\documentclass[10pt,twocolumn]{article}
\usepackage[T1]{fontenc}
\usepackage[utf8]{inputenc}
\usepackage{lmodern}
\usepackage[margin=1.8cm,top=2.4cm,bottom=2cm,columnsep=0.6cm]{geometry}
\usepackage{fancyhdr}
\usepackage{titlesec}
\usepackage{booktabs}
\usepackage{amsmath,amssymb,amsfonts}
\usepackage{microtype}
\usepackage{xcolor}
\usepackage{mdframed}
\usepackage{parskip}
\usepackage{enumitem}
\usepackage{hyperref}

\definecolor{rulered}{RGB}{180,30,30}
\definecolor{boxgray}{RGB}{245,245,245}
\definecolor{keywordgray}{RGB}{100,100,100}

\pagestyle{fancy}
\fancyhf{}
\fancyhead[L]{\textsc{\small Claude Mag}}
\fancyhead[C]{\small\textit{Machine Learning Research Digest}}
\fancyhead[R]{\small 2026-10-12}
\fancyfoot[C]{\small\thepage}
\renewcommand{\headrulewidth}{0.8pt}

\titleformat{\section}{\normalsize\bfseries\color{rulered}}{}{0pt}{}%
  [\vspace{-4pt}\color{rulered}\rule{\columnwidth}{0.4pt}]
\titlespacing{\section}{0pt}{8pt}{4pt}

\hypersetup{colorlinks=true,urlcolor=rulered,linkcolor=black}

\begin{document}
\setlength{\parindent}{0pt}
\setlength{\parskip}{4pt}

{\small\color{keywordgray}\textbf{Keywords:}
  mixture of experts \textbullet{}
  interpretability \textbullet{}
  sparse autoencoders \textbullet{}
  routing}\par\vspace{4pt}

{\LARGE\bfseries\raggedright
  RouterInterp: Understanding Superposed Specialisation in
  Mixture of Experts Routing}\par\vspace{4pt}

{\large\itshape\raggedright
  Experts Do Not Own Domains---They Activate on Disjoint Unions of
  Fine-Grained Features, and Sparse Autoencoder Probes Explain Why}\par\vspace{6pt}

{\small\textbf{By} Ilya Lasy, Nora Yinuo Cai \& Kola Ayonrinde}\par\vspace{2pt}

{\small\color{keywordgray}arXiv:2610.11775 \textbullet{} ICML 2026 \textbullet{} October 2026}
\par\vspace{2pt}\noindent{\color{rulered}\rule{\columnwidth}{0.6pt}}\vspace{6pt}

MoE routing does not follow the single-domain assumption:
each expert activates on a \emph{disjoint union} of fine-grained
features---the Superposed Specialisation Hypothesis. RouterInterp identifies
these features via Sparse Autoencoder probes and composes them into
natural-language routing explanations, achieving 65\% higher routing-detection
accuracy than token-statistics baselines on gpt-oss-20b.

\begin{mdframed}[backgroundcolor=boxgray,linecolor=rulered,linewidth=1pt,
    innerleftmargin=6pt,innerrightmargin=6pt,
    innertopmargin=6pt,innerbottommargin=6pt]
\textbf{\small AT A GLANCE}\par\vspace{4pt}
\begin{description}[leftmargin=0pt,labelsep=4pt,itemsep=2pt,parsep=0pt,
    font=\small\bfseries]
  \item[Problem:] It is unknown why specific tokens are routed to specific
    experts in MoE language models.
  \item[Why it matters:] Understanding routing is essential for auditing,
    debugging, and improving MoE architectures.
  \item[Method:] Sparse linear probes over SAE features predict routing;
    top-feature descriptions are composed into expert explanations.
  \item[Result:] $\approx$65\% higher routing-detection accuracy on
    gpt-oss-20b vs.\ token-statistics baselines.
  \item[Evidence:] Causal interventions confirm the SAE features causally
    influence routing; manual validation confirms explanation quality.
\end{description}
\end{mdframed}

\section{The Big Picture}

Mixture-of-Experts (MoE) models achieve state-of-the-art scaling
efficiency by routing each token through a small subset of expert
feed-forward networks. Despite widespread deployment, the routing
mechanism remains a black box: knowing \emph{why} a token is routed to
a particular expert is important for diagnosing failure modes, ensuring
balanced expert use, and building trust in deployed systems.

The dominant assumption---``the code expert,'' ``the French expert,'' etc.---is
a convenience fiction. RouterInterp provides empirical and causal evidence
for a richer picture: the \textbf{Superposed Specialisation Hypothesis (SSH)}.

\section{Why Existing Approaches Fall Short}

\textbf{Token-statistics methods} compute the frequency distribution of
token types per expert and label each with the most common category.
These methods are coarse, conflate co-occurrences with specialisation,
and show low prediction accuracy when routing is polysemous.

\textbf{Probing classifiers} on raw activations predict routing but
yield no human-readable explanation.

\textbf{Manual inspection} is intractable across the hundreds of experts
and billions of parameters in production MoE models.

\section{Core Mechanism}

\textbf{Superposed Specialisation Hypothesis:} Each expert $e$ activates
not on one broad domain, but on a disjoint union of fine-grained
linguistic or semantic features---mirroring the superposition phenomenon
known to occur in individual neurons.

RouterInterp operationalises this via Sparse Autoencoders (SAEs):
\begin{enumerate}[itemsep=2pt,parsep=0pt]
  \item Derive monosemantic SAE features from the residual stream.
  \item Fit $L_1$-penalised probes to predict routing from SAE activations.
  \item Compose the top-predictive features' descriptions into a
    unified expert explanation.
\end{enumerate}

\section{Technical Formulation}

Let $\mathbf{h}_t \in \mathbb{R}^d$ be the residual stream at token $t$.
A Sparse Autoencoder maps it to a sparse feature vector:
\[
\mathbf{f}_t = \mathrm{ReLU}(\mathbf{W}_\text{enc}\,\mathbf{h}_t + \mathbf{b}_\text{enc}),
\quad \mathbf{f}_t \in \mathbb{R}^F
\]

The routing indicator for expert $e$ is:
\[
r_{t,e} = \mathbf{1}[\text{expert } e \text{ selected for token } t]
\]

RouterInterp solves the $L_1$-penalised logistic regression:
\[
\hat{r}_{t,e} = \sigma(\boldsymbol{\beta}_e^\top \mathbf{f}_t),
\quad \boldsymbol{\beta}_e = \arg\min_{\boldsymbol{\beta}}
\mathcal{L}_\text{CE}(r,\hat{r}) + \mu\|\boldsymbol{\beta}\|_1
\]

Non-zero entries $\{k : \beta_{e,k} \neq 0\}$ identify the SAE features
predictive of routing to expert $e$. Their descriptions $\{s_k\}$ are
composed into a natural-language explanation via a language model.

\section{Learning or Inference Procedure}

\begin{enumerate}[itemsep=2pt,parsep=0pt]
  \item Train a SAE on the frozen MoE's residual stream to obtain a
    dictionary of monosemantic features.
  \item Collect routing labels $r_{t,e}$ over a large corpus.
  \item Fit sparse probes $\boldsymbol{\beta}_e$ for each expert over SAE
    feature activations.
  \item Extract the top-$k$ predictive features per expert.
  \item Pass the feature descriptions to a language model to compose
    a coherent expert summary.
\end{enumerate}

\section{What the Guarantee Says}

RouterInterp provides an empirical claim validated by causal intervention:
activating a high-weight SAE feature at a token that would otherwise not
route to expert $e$ measurably increases routing probability to $e$,
confirming causality beyond mere correlation.

\section{Experimental Findings}

\begin{itemize}[itemsep=2pt,parsep=0pt]
  \item On \textbf{gpt-oss-20b}, RouterInterp achieves $\approx$\textbf{65\%
    higher routing-detection accuracy} than token-statistics baselines.
  \item Manual inspection confirms that expert descriptions are specific,
    coherent, and capture non-obvious expert roles.
  \item The SSH is validated: the top SAE features per expert span
    semantically disjoint categories (e.g., medical terminology
    \emph{and} certain syntactic patterns in Spanish for the same expert).
  \item Results are consistent across all expert layers in gpt-oss-20b.
\end{itemize}

\section{Ablations and Interpretation}

\textbf{Raw activations vs.\ SAE features:} Probing on raw residual stream
activations degrades routing prediction accuracy significantly and yields
incoherent explanations.

\textbf{Top-$k$ features:} Including more features beyond the natural
coefficient elbow adds noise without improving routing accuracy.

\textbf{Causal intervention:} Artificially activating a predictive SAE feature
causally increases routing probability to the predicted expert.

\textbf{Across layers:} Early layers show routing correlated with syntactic
SAE features; later layers show richer semantic specialisation.

\section*{Reference}

Ilya Lasy, Nora Yinuo Cai, Kola Ayonrinde.
\textbf{RouterInterp: Understanding Superposed Specialisation in Mixture of
Experts Routing}.
ICML 2026. arXiv:2610.11775, October 2026.\\
\url{https://arxiv.org/abs/2610.11775}

\end{document}
